\begin{abstract}
Emerging studies, exemplified by BitNet~\cite{bitnet}, are ushering in a transformative phase characterized by 1-bit Large Language Models (LLMs). This paper presents a ternary 1-bit LLM variant, \textbf{\text{BitNet b1.58}}, where each parameter (or weight) in the network is drawn from the set \{-1, 0, 1\}. BitNet b1.58 attains parity with full-precision Transformer LLMs (utilizing FP16 or BF16) of equivalent scale and trained on identical corpora, achieving similar results with respect to perplexity and downstream tasks. Crucially, it yields notable gains in latency, memory consumption, computational throughput, and energy efficiency. Perhaps even more significantly, the 1.58-bit LLM establishes an updated scaling principle and training methodology, thereby guiding the creation of high-performing yet resource-efficient LLMs for future generations. Additionally, this development heralds novel avenues for algorithm design and hardware that are purpose-built for 1-bit LLM computations.
\end{abstract}

\section{The Era of 1-bit LLMs}

Artificial Intelligence has witnessed an unprecedented surge in the scale and sophistication of Large Language Models (LLMs) over recent years. While these models have set new benchmarks in various natural language processing benchmarks, their burgeoning parameter counts present obstacles to real-world deployment and exacerbate concerns regarding environmental sustainability and operational cost due to escalating energy demands.

One promising strategy to alleviate these issues involves post-training quantization, where models are compressed to lower numerical precisions for inference~\cite{smoothquant, gptq, quip, quip_sharp}. Such quantization dramatically lessens the burden on memory and compute by reducing the bit-width of both weights and activations. There has been a clear trajectory toward decreasing the precision from traditional 16-bit formats to even 4-bit alternatives~\cite{gptq, awq}. However, despite its widespread use across industry-grade LLM deployments, post-training quantization is inherently sub-optimal and leaves performance on the table.

Recent advances in 1-bit model architectures, such as BitNet~\cite{bitnet}, have shown considerable promise in lowering the computational expense of large language models (LLMs) while preserving their effectiveness. Standard LLMs utilize 16-bit floating-point representations (such as FP16 or BF16), and their computations are dominated by matrix multiplications. Consequently, the primary source of computational overhead is the frequent execution of floating-point multiplications and additions. BitNet, in contrast, performs matrix multiplications using only integer additions, which substantially reduces energy consumption in LLMs. Given that power is often the bottleneck in modern hardware accelerators, these energy reductions can translate directly into accelerated processing speeds.

Beyond computational considerations, transferring model weights from off-chip DRAM to the on-chip accelerator’s memory (such as SRAM) is another major bottleneck during inference. While increasing on-chip SRAM can raise throughput, the resulting costs can far exceed those of DRAM. 1-bit LLMs stand out from their full-precision counterparts by requiring significantly less memory both in terms of storage capacity and bandwidth. As a result, they enable much faster and less costly weight transfers from DRAM, thus enhancing inference speed and efficiency.

In this study, we propose a new 1-bit LLM variant named \textbf{\text{BitNet b1.58}}, where all parameters are ternary, restricted to the set \{-1, 0, 1\}. By augmenting the original 1-bit BitNet design with a zero value, the effective parameter precision reaches 1.58 bits per parameter. \text{BitNet b1.58} retains all key properties of its 1-bit predecessor, notably a computational framework that relies almost entirely on additions, eliminating most multiplication operations and allowing for highly efficient implementations. It shares the energy efficiency of the original BitNet while offering major improvements over FP16-based models in terms of memory usage, data throughput, and latency. Additionally, \text{BitNet b1.58} introduces two notable enhancements: first, it offers stronger representational capacity by explicitly supporting feature suppression—thanks to the presence of zero weights—thereby boosting the performance potential of 1-bit LLMs; and second, our empirical results demonstrate that \text{BitNet b1.58} achieves perplexity and end-task results comparable to full-precision (FP16) baselines, beginning at the 3B model scale, provided identical training regimes and configurations are used.

\section{BitNet b1.58}

\text{BitNet b1.58} builds upon the BitNet framework, a Transformer variant where the conventional \emph{nn.Linear} layers are substituted with \emph{BitLinear}. The model is initialized and trained from the ground up using weights quantized to 1.58 bits and activations quantized to 8 bits. This version incorporates several changes compared to the original BitNet, which are outlined as follows.

\paragraph{Quantization Function.} We employ an \emph{absmean} quantization function to restrict weights to values in the set $\{-1, 0, +1\}$. The process involves normalizing the weight matrix by its mean absolute value, after which each element is mapped to the nearest integer in $\{-1, 0, +1\}$ through rounding and clipping:

\begin{equation}
    \widetilde{W} = \mathrm{RoundClip}(\frac{W}{\gamma + \epsilon}, -1, 1),
\end{equation}

\begin{equation}
    \mathrm{RoundClip}(x, a, b) = \max(a, \min(b, \mathrm{round}(x))),
\end{equation}

\begin{equation}
    \gamma = \frac{1}{nm}\sum_{ij} |W_{ij}|.
\end{equation}

For the activations, the quantization approach remains consistent with that of BitNet, except for the scaling strategy. In contrast to scaling pre-activation values into $[0, Q_b]$ before applying non-linear functions, activations in our model are uniformly scaled to $[-Q_b, Q_b]$ for each token, eliminating the need for zero-point quantization. This alternative method simplifies both coding and system-level optimization without producing any substantial difference in empirical performance.

\paragraph{LLaMA-alike Components.}

The LLaMA~\cite{llama, llama2} design has become a standard foundation for open-source large language models. To align with this widely adopted structure, \text{BitNet b1.58} incorporates LLaMA-inspired modules, specifically integrating RMSNorm~\cite{rmsnorm}, SwiGLU~\cite{swiglu}, rotary embedding~\cite{rope}, and omitting all bias terms. These choices ensure that \text{BitNet b1.58} can be readily utilized with common open-source tools and libraries (such as Huggingface, vLLM~\cite{vllm}, and llama.cpp\footnote{\url{https://github.com/ggerganov/llama.cpp}}) with minimal modification required.

\section{Results}

We benchmarked \text{BitNet b1.58} against our own FP16 implementation of the \text{LLaMA LLM} across a range of model sizes. For a controlled evaluation, both models were pre-trained with 100 billion tokens sourced from the RedPajama dataset~\cite{redpajama}. Zero-shot performance was assessed on diverse language benchmarks, including ARC-Easy~\cite{arc}, ARC-Challenge~\cite{arc}, Hellaswag~\cite{hellaswag}, Winogrande~\cite{winoGrande}, PIQA~\cite{piqa}, OpenbookQA~\cite{openbookqa}, and BoolQ~\cite{boolq}. In addition, validation perplexities were measured using the WikiText2~\cite{wikitext2} and C4~\cite{c4} corpora.

For profiling memory consumption and inference speed, we examined runtime GPU usage and latency for both \text{LLaMA LLM} and \text{BitNet b1.58}. All timing was obtained with the FasterTransformer\footnote{\url{https://github.com/NVIDIA/FasterTransformer}} framework, which provides highly optimized large language model inference on GPUs. In these experiments, we incorporated the 2-bit kernel from Ladder~\cite{ladder} when running \text{BitNet b1.58}. The primary metric was per-token generation latency, reflecting the main inference cost.

Table~\ref{tab:ppl} provides a comparison of perplexity and resource demands for \text{BitNet b1.58} and \text{LLaMA LLM}. Notably, \text{BitNet b1.58} achieves comparable perplexity to full-precision \text{LLaMA LLM} from the 3B parameter scale onward, with a speedup of 2.71x and a 3.55x reduction in GPU memory footprint. Specifically, at 3.9B parameters, \text{BitNet b1.58} delivers 2.4x faster inference, uses 3.32x less memory, and outperforms the 3B \text{LLaMA LLM}.

A detailed breakdown of zero-shot accuracy on downstream tasks is shown in Table~\ref{tab:zero-shot}. The evaluation procedure followed the \emph{lm-evaluation-harness}\footnote{\url{https://github.com/EleutherAI/lm-evaluation-harness}} toolkit. Our results indicate that as the model size grows, the difference in accuracy between \text{BitNet b1.58} and \text{LLaMA LLM} diminishes. Of particular note, \text{BitNet b1.58} matches full-precision baseline accuracy at the 3B model size. Consistent with the perplexity analysis, the larger 3.9B \text{BitNet b1.58} outperforms the 3B \text{LLaMA LLM} on downstream tasks, while retaining advantages in memory and speed. Overall, these findings illustrate that \text{BitNet b1.58} represents a Pareto improvement over leading LLM baselines.

\paragraph{Memory and Latency}

We extended our experiments to larger models, including those with 7B, 13B, and 70B parameters, and analyzed the corresponding costs. As depicted in Figure~\ref{fig:latency-memory-energy}, latency and memory usage both display clear trends with model scaling: larger models benefit from greater speed improvements. Notably, \text{BitNet b1.58} 70B demonstrates a 4.1-fold speed advantage over the \text{LLaMA LLM} reference. This can be attributed to the increasing computational overhead of the \emph{nn.Linear} operation as models become larger. Memory requirements show a similar scaling pattern, primarily because embeddings are kept in full precision, making up a decreasing share of the overall memory footprint as models grow. All latency and memory metrics were collected using a 2-bit kernel; thus, further efficiency gains are possible with kernel-level optimization.

\paragraph{Energy}

Our evaluation also includes an estimation of the energy consumed by arithmetic operations for both \text{BitNet b1.58} and \text{LLaMA LLM}. Our primary focus is on matrix multiplication, as it constitutes the bulk of computational expenditure in LLMs. Figure~\ref{fig:energy} provides a breakdown of energy usage, revealing that \text{BitNet b1.58} primarily relies on INT8 addition, whereas \text{LLaMA LLM} incorporates both FP16 addition and FP16 multiplication. Based on the energy consumption models from~\cite{energycost, pokebnn}, \text{BitNet b1.58} achieves a 71.4$\times$ reduction in energy use for matrix multiplication when running on 7nm hardware. We also measured the total energy required for processing sequences of 512 tokens, observing that \text{BitNet b1.58}'s energy efficiency advantage over the FP16 \text{LLaMA LLM} increases with model size. This outcome stems from the greater relative contribution of \emph{nn.Linear} operations at larger scales, while other model components account for a smaller fraction of the total energy demand in bigger models.

\paragraph{Throughput}

We assess and compare the throughput of \text{BitNet b1.58} and the 70B parameter \text{LLaMA LLM} on two 80GB A100 GPUs, employing pipeline parallelism~\cite{gpipe} to make the execution of the \text{LLaMA LLM} 70B feasible. The batch size was incrementally increased until reaching the GPU memory constraints, with all tests using a sequence length of 512. According to Table~\ref{tab:throughput}, the 70B \text{BitNet b1.58} model can handle a batch size up to 11 times larger than \text{LLaMA LLM}, resulting in an 8.9-fold improvement in throughput.

\textbf{\text{BitNet b1.58} introduces a novel scaling law that relates model performance and inference efficiency}. To illustrate, based on Figures~\ref{fig:latency-memory-energy} and \ref{fig:energy}, we present the following comparisons of equivalency between different model capacities using 1.58-bit and 16-bit precision:

\begin{itemize}
    \item BitNet b1.58 with 13B parameters surpasses 3B FP16 LLM in latency, memory footprint, and energy efficiency.
    \item BitNet b1.58 with 30B parameters outperforms 7B FP16 LLM regarding latency, memory requirements, and energy usage.
    \item BitNet b1.58 at 70B parameters exceeds the efficiency of 13B FP16 LLM across latency, memory consumption, and energy.
\end{itemize}

\paragraph{Training with 2T Tokens}

The volume of training tokens is an essential determinant in the performance of LLMs. To evaluate the token-level scalability of \text{BitNet b1.58}, we conducted training of a \text{BitNet b1.58} model on 2T tokens, utilizing the data strategy from StableLM-3B~\cite{StableLM-3B-4E1T}, which stands as a state-of-the-art 3B open-source model. The performance was measured on a benchmark suite consisting of Winogrande~\cite{winoGrande}, PIQA~\cite{piqa}, SciQ~\cite{sciq}, LAMBADA~\cite{lambada}, and ARC-easy~\cite{arc}, with zero-shot accuracy results documented in Table~\ref{tab:2t}. For tasks assessed by both accuracy and normalized accuracy, we report their average. The StableLM 3B results at 2T tokens are quoted from its original technical documentation. Our analysis demonstrates that \text{BitNet b1.58} attains higher performance across all evaluated tasks, suggesting that LLMs utilizing 1.58-bit quantization possess robust generalization properties.

\section{Discussion and Future Work}

\textbf{1-bit Mixture-of-Experts (MoE) LLMs}

Mixture-of-Experts (MoE) architectures have established themselves as a resource-efficient technique for large language models (LLMs), offering substantial reductions in computational FLOPs. However, high memory demands and inter-chip communication overhead currently constrain their scalability and adoption. The advent of 1.58-bit LLMs provides a potential solution to these bottlenecks. With a lower memory footprint, fewer devices are needed to support MoE models, and the associated communication overhead for transmitting activations across devices is significantly decreased. Ideally, if the entire MoE model fits onto a single chip, communication overhead could be eliminated altogether.

\textbf{Native Support of Long Sequence in LLMs}

The capability to process long sequences has become increasingly crucial for modern LLMs. A principal challenge for long sequence inference is the high memory usage caused by KV cache storage. \text{BitNet b1.58} offers important progress toward native support for extended sequence lengths by compressing activations from 16 bits to 8 bits, thus permitting the context window to double with fixed hardware resources. Further, for 1.58-bit LLMs, this activation memory can be losslessly reduced to as low as 4 bits or possibly even less—an avenue we propose for subsequent exploration.

\textbf{LLMs on Edge and Mobile}

Deploying 1.58-bit LLMs on edge and mobile hardware stands to markedly elevate their practical usability. Such devices are often bounded by stringent constraints on memory and compute capacity, which in turn restrict the feasible scale of LLMs. However, the efficiency gains of 1.58-bit LLMs—both in terms of memory footprint and energy use—allow for practical deployment on edge and mobile platforms, thus unlocking a diverse set of novel applications. This innovation could greatly expand the operational scope and power of these devices. Furthermore, because 1.58-bit LLMs are inherently better aligned with CPU architectures—prevalent in edge and mobile settings—\text{BitNet b1.58} can run with enhanced efficiency, boosting performance and broadening real-world applicability.

\textbf{New Hardware for 1-bit LLMs}

Initiatives such as Groq\footnote{\url{https://groq.com/}} have begun illustrating the advantages of designing dedicated hardware, such as language processing units (LPUs), for LLM workloads. Looking ahead, we advocate for focused research and development into new hardware and computing systems tailored specifically for 1-bit LLMs, motivated by the distinct computational paradigm that BitNet~\cite{bitnet} has introduced.

\end{document}